% Part: first-order-logic
% Chapter: natural-deduction
% Section: provability-propositional

% verification of properties of provability needed for maximally
% consistent sets in the completeness chapter.

\documentclass[../../../include/open-logic-section]{subfiles}

\begin{document}

\iftag{FOL}
      {\olfileid[id]{fol}{ntd}{ppr}}
      {\olfileid[id]{pl}{ntd}{ppr}}

\olsection{\usetoken{S}{derivability} dan Penghubung Proposisional}

\begin{explain}
  Kita membuktikan bahwa relasi !!{derivability}~$\Proves$ dari deduksi
  alami cukup kuat untuk membuktikan beberapa fakta dasar yang melibatkan
  penghubung proposisional, seperti $!A \land !B \Proves !A$ dan
  $!A, !A \lif !B \Proves !B$ (modus ponens). Fakta-fakta ini diperlukan
  dalam pembuktian teorema kelengkapan.
\end{explain}

\begin{prop}\ollabel{prop:provability-land}
  \begin{enumerate}
  \item \ollabel{prop:provability-land-left} Baik $!A \land !B \Proves
    !A$ maupun $!A \land !B \Proves !B$.
  \item \ollabel{prop:provability-land-right} $!A, !B \Proves !A \land !B$.
  \end{enumerate}
\end{prop}

\begin{proof}
  \begin{enumerate}
  \item Kita dapat !!{derive} keduanya:
    \begin{prooftree}
      \AxiomC{$!A \land !B$}
      \RightLabel{\Elim{\land}}
      \UnaryInfC{$!A$}
      \DisplayProof\qquad\bottomAlignProof
      \AxiomC{$!A \land !B$}
      \RightLabel{\Elim{\land}}
      \UnaryInfC{$!B$}
    \end{prooftree}
  \item Kita dapat !!{derive}:
    \begin{prooftree}
      \AxiomC{$!A$}
      \AxiomC{$!B$}
      \RightLabel{\Intro{\land}}
      \BinaryInfC{$!A \land !B$}
    \end{prooftree}
  \end{enumerate}
\end{proof}


\begin{prop}\ollabel{prop:provability-lor}
  \begin{enumerate}
  \item $!A \lor !B, \lnot !A, \lnot !B$ tak konsisten.
  \item Baik $!A \Proves !A \lor !B$ maupun $!B \Proves !A \lor !B$.
  \end{enumerate}
\end{prop}

\begin{proof}
  \begin{enumerate}
  \item Perhatikan !!{derivation} berikut:
    \begin{prooftree}
      \AxiomC{$!A \lor !B$}
      \AxiomC{$\lnot !A$}
      \AxiomC{$\Discharge{!A}{1}$}
      \RightLabel{\Elim{\lnot}}
      \BinaryInfC{$\lfalse$}
      \AxiomC{$\lnot !B$}
      \AxiomC{$\Discharge{!B}{1}$}
      \RightLabel{\Elim{\lnot}}
      \BinaryInfC{$\lfalse$}
      \DischargeRule{\Elim{\lor}}{1}
      \TrinaryInfC{$\lfalse$}
    \end{prooftree}
    Ini merupakan !!a{derivation} untuk~$\lfalse$ dari asumsi
    !!{undischarged} $!A \lor !B$, $\lnot !A$, dan $\lnot !B$.
  \item Kita dapat !!{derive} keduanya:
    \begin{prooftree}
      \AxiomC{$!A$}
      \RightLabel{\Intro{\lor}}
      \UnaryInfC{$!A \lor !B$}
      \DisplayProof\qquad\bottomAlignProof
      \AxiomC{$!B$}
      \RightLabel{\Intro{\lor}}
      \UnaryInfC{$!A \lor !B$}
    \end{prooftree}
  \end{enumerate}
\end{proof}

\begin{prop}\ollabel{prop:provability-lif}
  \begin{enumerate}
  \item \ollabel{prop:provability-lif-left} $!A, !A \lif !B \Proves !B$.
  \item \ollabel{prop:provability-lif-right}
    Baik $\lnot !A \Proves !A \lif !B$ maupun $!B \Proves !A \lif !B$.
  \end{enumerate}
\end{prop}

\begin{proof}
  \begin{enumerate}
  \item Kita dapat !!{derive}:
    \begin{prooftree}
      \AxiomC{$!A \lif !B$}
      \AxiomC{$!A$}
      \RightLabel{\Elim{\lif}}
      \BinaryInfC{$!B$}
    \end{prooftree}

  \item Hal ini ditunjukkan oleh dua !!{derivation} berikut:
    \begin{prooftree}
      \AxiomC{$\lnot !A$}
      \AxiomC{$\Discharge{!A}{1}$}
      \RightLabel{\Elim{\lnot}}
      \BinaryInfC{$\lfalse$}
      \RightLabel{\FalseInt}
      \UnaryInfC{$!B$}
      \DischargeRule{\Intro{\lif}}{1}
      \UnaryInfC{$!A \lif !B$}
      \DisplayProof\qquad\bottomAlignProof
      \AxiomC{$!B$}
      \RightLabel{\Intro{\lif}}
      \UnaryInfC{$!A \lif !B$}
    \end{prooftree}
    Perhatikan bahwa $\Intro{\lif}$ dapat !!{discharge} nol atau lebih
    kemunculan asumsi~$!A$; dalam !!{derivation} di sebelah kanan, aturan
    itu diterapkan tanpa adanya asumsi yang dilepaskan.
  \end{enumerate}
\end{proof}

\end{document}
